\section{A cyclotomic annihilator for the signed Selmer group}
\label{sec:selmer-annihilator}

Let \(N\ge1\), let \(E\subset\mathbb Q(\mu_N)\) be a CM field, and put
\(d=[E:\mathbb Q]\), \(\Gamma=\operatorname{Gal}(E/\mathbb Q)\).
Write \(c\) for complex conjugation in \(\Gamma\). Fix a prime
\(\ell\) and an integer \(e\ge1\). We assume either
\(\ell\ne2\), or \(\ell=2\) and \(\mu_4\subset E\).
Put \(k=e+1\), \(\epsilon=(-1)^k\), and
\(A=(\mathbb Q_\ell/\mathbb Z_\ell)(-e)\). The action on \(A\)
is multiplication by \(\chi_\ell^{-e}\), where \(\chi_\ell\) is
the cyclotomic character. Let \(S\) be the set of places of \(E\)
over \(\ell\). With \(G_{E,S}\) the quotient of the absolute
Galois group unramified outside \(S\) at finite places, set
\[
 H=H^1_{\rm cts}(G_{E,S},A),\qquad
 P=\ker\!\left(H\longrightarrow\prod_{v\in S}H^1(E_v,A)\right),
 \qquad P_\epsilon=\{x\in P:cx=\epsilon x\}.
\]
The action of \(\Gamma\) on \(H\) uses both conjugation on the
Galois group and the given cyclotomic action on the coefficients.
These groups are the \(H^1\), \(\mathrm{Sha}\) and
\(\mathrm{ShaMinus}\) of the source for these data.

\paragraph{Annihilator and length bounds.}
There is an integer \(r\ge0\), independent of \(x\in P_\epsilon\),
such that
\begin{equation}
 \ell^r P_\epsilon=0,\qquad
 r\log\ell\le 30d(e+1)\log(e\ell N).
 \label{eq:cyclotomic-kill}
\end{equation}
Consequently \(P_\epsilon\) has finite \(\mathbb Z_\ell\)-length and
\begin{equation}
 \operatorname{length}_{\mathbb Z_\ell}(P_\epsilon)\log\ell
 \le 240\ell d^2(e+1)\log^2(e\ell N).
 \label{eq:cyclotomic-length}
\end{equation}
In particular these inequalities apply at odd primes dividing \(d\).
The separate characterwise estimate in the prime-to-\(d\) case is a
stronger assertion.

\subsection{The weight-one class-group relation}

For \(L>2\) put \(K_L=\mathbb Q(\mu_L)\),
\(G_L=(\mathbb Z/L\mathbb Z)^\times\), and let \(\sigma_b\) send
\(\zeta_L\) to \(\zeta_L^b\). In formulas at level \(L\), the
symbol \([b]_L\) denotes the representative in \([0,L-1]\).
Fix an odd positive integer \(a\) prime to \(L\), and write
\[
 j_L(b)=\frac{a[a^{-1}b]_L-[b]_L}{L},\qquad
 h_L(b)=\frac{a-1}{2}-j_L(b).
\]
Both numbers are integers. For the usual Stickelberger element
\(\vartheta_L=L^{-1}\sum_{b\in G_L}[b]_L\sigma_b^{-1}\), coefficient
comparison gives the identity in \(\mathbb Q[G_L]\)
\[
 \sum_{b\in G_L}j_L(b)\sigma_b^{-1}
   =(a\sigma_a^{-1}-1)\vartheta_L.
\]
The right side belongs to \(\mathbb Z[G_L]\vartheta_L\cap
\mathbb Z[G_L]\). Stickelberger's theorem says that this ideal
annihilates the ideal class group of \(K_L\), for every cyclotomic
modulus \(L\). Also, \(\sum_b\sigma_b^{-1}\) annihilates that class
group, because the product of the conjugates of a fractional ideal
is the extension of its norm to \(\mathbb Q\), a principal ideal.
Subtracting the two relations gives
\begin{equation}
 \sum_{b\in G_L}h_L(b)\sigma_b^{-1}C=0
 \quad\text{for every } C\in\operatorname{Cl}(K_L)
 \label{eq:cyclotomic-class}
\end{equation}
in additive notation. The precise classical input, with these group
ring conventions, is R.~Sharifi, \emph{Iwasawa Theory}, Definition
4.5.1 and Theorem~4.5.6. The same theorem is proved in
L.~Washington, \emph{Introduction to Cyclotomic Fields}, second
edition, Theorem~6.10.

Suppose \(\eta:G_{K_L}\to A\) is a continuous homomorphism with
finite image, written with trivial coefficient action, unramified
at every finite place. The field \(K_L\) is totally imaginary.
The Hilbert class field theorem therefore identifies the finite
unramified abelian quotient through which \(\eta\) factors with a
quotient of \(\operatorname{Cl}(K_L)\). Artin reciprocity is equivariant
for the action of \(G_L\). For \(u\in G_{K_L}\) and any absolute
lifts \(\widetilde\sigma_b\) of \(\sigma_b\), applying \(\eta\) to
the Artin image of \eqref{eq:cyclotomic-class} gives
\begin{equation}
 \sum_b h_L(b)\eta(\widetilde\sigma_b^{-1}u\widetilde\sigma_b)=0.
 \label{eq:cyclotomic-hom}
\end{equation}
Conjugation with the inverse lift on the left gives the inverse
automorphism on the ideal class, as required by \eqref{eq:cyclotomic-class}.

\subsection{Lifting the weight and passing to a fixed group ring}

Set \(M=\ell N\), \(a=2M+1\), and let \(L=M\ell^m\) with \(m\ge0\).
Write \(\pi_L:G_L\to\Gamma\) for restriction. The ordinary Bernoulli
polynomials have \(B_1(X)=X-\tfrac12\). Define
\begin{align*}
 E_k(L,x)&=\frac{L^{k-1}}{k} B_k(x/L),\\
 \Theta(L)&=-\sum_{b\in G_L} E_k(L,[b]_L)\pi_L(b),\\
 Y(L)&=(1-a^k\pi_L(a))\Theta(L),\\
 W(L)&=-\sum_{b\in G_L}[b]_L^{\,e}h_L(b)\pi_L(b).
\end{align*}
Here \(\Theta(L),Y(L)\in\mathbb Q[\Gamma]\) and
\(W(L)\in\mathbb Z[\Gamma]\).
The source evaluates its inverse-indexed group-ring sum on the
inverse of the restriction map. This accounts for \(\pi_L(b)\)
as written in these three expressions.

The two needed identities are
\[
 Y(L)=Y(M),\qquad
 Y(M)-W(L)\in\ell^t\mathbb Z_\ell[\Gamma]
 \quad\bigl(m\ge t+v_\ell(D_k)+1,\ t\ge0\bigr),
 \qquad
 D_k=2k\prod_{p\le k+1}p ,
\]
where the product is over primes. For the equality, every lift of
a unit modulo \(M\) to a residue modulo \(M\ell^m\) is a unit,
since \(\ell\mid M\). Summing \(E_k(L,x)\) over its \(\ell^m\)
representatives gives \(E_k(M,[x]_M)\) by
\[
 \sum_{i=0}^{q-1}B_k\!\left(\frac{u+i}{q}\right)
    =q^{1-k}B_k(u),\qquad q=\ell^m .
\]
This distribution identity follows by summing the generating series
\(te^{ut}/(e^t-1)\). Restriction to \(\Gamma\) is constant on each
fiber, proving the assertion for \(\Theta\) and hence for \(Y\).

For the congruence, put \(P_k=\prod_{p\le k+1}p\).
The binomial expansion of \(B_k(X)\) and the von Staudt--Clausen
theorem imply
\[
 E_k(L,x)=\frac{x^k}{kL}-\frac{x^{k-1}}2+
               \frac{L}{kP_k}z_x,\qquad z_x\in\mathbb Z .
\]
Indeed, the term indexed by \(i\ge2\) in the Bernoulli expansion has
a factor \(L^{i-1}\), and \(P_k B_i\) is integral for \(2\le i\le k\).
Set \(x=[b]_L\), \(y=[a^{-1}b]_L\), \(j=j_L(b)\), so that
\(ay=x+Lj\). For some integers \(u,v\),
\[
 (ay)^k=x^k+kx^{k-1}Lj+L^2u,\qquad
 (ay)^{k-1}=x^{k-1}+Lv.
\]
Substitution yields
\[
 E_k(L,x)-a^kE_k(L,y)-x^{k-1}h_L(b)
 =\frac{L}{D_k}
    \bigl(2z_x-2a^k z_y-2P_k u+akP_k v\bigr).
\]
Multiplying by \(-\pi_L(b)\), summing and using \(e=k-1\) proves
the congruence. Taking \(t=0\) also gives \(Y:=Y(M)\in
\mathbb Z_\ell[\Gamma]\). Thus a single rational group-ring element
is integral at \(\ell\) and is approximated by the explicit \(W(L)\)
to every required precision.

\subsection{The Selmer restriction and cancellation of the twist}

Let \(U=G_{E(\mu_{\ell^\infty})}\). If \(x\in P\), choose a
continuous cocycle \(f:G_E\to A\) inflated from \(G_{E,S}\)
representing \(x\). It has finite image by compactness of \(G_E\)
and discreteness of \(A\). For each \(v\mid\ell\), choose a
decomposition group \(D_v\) in \(G_E\). Local triviality provides
\(t_v\in A\) such that
\(f(g)=(\chi_\ell(g)^{-e}-1)t_v\) on \(D_v\). Choose \(s\ge1\)
so that \(\ell^s\) kills the image of \(f\) and all the finitely
many \(t_v\). On \(U\), \(f\) is a homomorphism \(\kappa\);
coboundaries vanish there. In particular \(\kappa\) depends only
on \(x\).

For every \(m\ge s\), \(K_L=\mathbb Q(\mu_{M\ell^m})\) contains
\(E\) and \(\mu_{\ell^s}\). Put \(V_L=G_{K_L}\).
The restriction \(\eta=f|_{V_L}\) is a homomorphism, since \(V_L\)
acts trivially on the image of \(f\). We show it is unramified.
Away from \(\ell\), the relevant inertia is contained in inertia of
\(E\) killed by \(G_E\to G_{E,S}\), so \(f=0\) there. At \(v\mid\ell\)
the character of \(V_L\cap D_v\) acts trivially on \(t_v\), giving
\(f|_{V_L\cap D_v}=0\). For any \(h\in G_E\) and \(w\in V_L\), the
cocycle rule gives
\[
 f(hwh^{-1})=\chi_\ell(h)^{-e} f(w)
                    +(1-\chi_\ell(w)^{-e})f(h).
\]
The last term vanishes since \(\ell^s f(h)=0\). Conjugating the
chosen decomposition groups therefore covers every place above
\(\ell\) in \(K_L\). This proves the claimed ramification property
of \(\eta\) on the entire finite layer.

If \(u\in U\cap G_{K_M}\), it fixes all \(M\ell^m\)-th roots of
unity. Indeed \(u\) fixes \(\mu_M\) and every \(\ell\)-power root
of unity, and these roots generate \(\mu_{M\ell^m}\).
For \(b\in G_L\) put \(w_b=\widetilde\sigma_b^{-1}u\widetilde\sigma_b\).
Both \(V_L\) and \(U\) are normal under these lifts, so \(w_b\) lies
in their intersection. Thus \(\kappa(w_b)=\eta(w_b)\).
The action on cohomology, evaluated on \(u\), is
\[
 (\pi_L(b)x)|_U(u)=
             \chi_\ell(\widetilde\sigma_b)^{-e}\,\kappa(w_b).
\]
Changing the lift changes the representing cocycle by a coboundary,
which vanishes on \(U\). Also,
\(\chi_\ell(\widetilde\sigma_b)\equiv[b]_L\bmod\ell^s\). Hence
\([b]_L^e\chi_\ell(\widetilde\sigma_b)^{-e}\) acts as the identity on
\(\kappa(w_b)\). Equation~\eqref{eq:cyclotomic-hom} now gives
\[
 (W(L)x)|_U(u)=-\sum_b h_L(b)\eta(w_b)=0 .
\]
Taking \(m\ge s+v_\ell(D_k)+1\) gives \(W(L)x=Yx\), since
\(\ell^s x=0\). We have proved that the homomorphism \((Yx)|_U\)
vanishes on \(U\cap G_{K_M}\). For every \(u\in U\), the power
\(u^{\varphi(M)}\) fixes \(\mu_M\). Thus
\((\varphi(M)Yx)|_U=0\).

We justify the final passage back to global cohomology. Inflation
from \(G_{E,S}\) to \(G_E\) is injective on \(H^1\) because the
coefficient action factors through \(G_{E,S}\). A cocycle on \(G_E\)
whose restriction vanishes on \(U\) descends to
\(J=G_E/U=\chi_\ell(G_E)\) by the cocycle rule. This is an open
subgroup of \(\mathbb Z_\ell^\times\). Indeed its index equals
\([E\cap\mathbb Q(\mu_{\ell^\infty}):\mathbb Q]\), which is finite.
For odd \(\ell\), the group of units is the product of a finite
cyclic group of order prime to \(\ell\) and \(1+\ell\mathbb Z_\ell
\simeq\mathbb Z_\ell\), hence is procyclic. Its closed subgroups
are procyclic. For \(\ell=2\) our hypothesis puts \(J\) in
\(1+4\mathbb Z_2\simeq\mathbb Z_2\). In both cases an element
\(\gamma\) topologically generating \(J\) has infinite order.
The nonzero scalar \(\gamma^{-e}-1\) acts surjectively on the
divisible group \(A\). Given a continuous cocycle \(z\) on \(J\),
choose \(t\) with \(z(\gamma)=(\gamma^{-e}-1)t\). Subtract this
coboundary. The resulting cocycle vanishes on all integral powers
of \(\gamma\), and then on \(J\) by continuity. This proves
\[
             \varphi(M)Y P=0 .
\]

\subsection{The determinant on the signed quotient}

Let \(R=\mathbb Z_\ell[\Gamma]/(c-\epsilon)\). This is a free
\(\mathbb Z_\ell\)-module of rank \(d/2\). In fact \(c\) has order
two since \(E\) is CM and each pair \(g,cg\) contributes the
single basis element represented by \(g\), with \(cg=\epsilon g\).
The same integral basis works over \(\mathbb Z,\mathbb Q,\mathbb C\).
Write \(\Delta\) for the determinant of multiplication by \(Y\) in
this basis. For \(x\in P_\epsilon\), the map \(R\to H\) sending
\(b\mapsto b\varphi(M)x\) is well-defined. The image of \(YR\)
under this map vanishes. By the adjugate identity,
\(\Delta\cdot1\) belongs to \(YR\), whence
\begin{equation}
                  \Delta\varphi(M)x=0 .
 \label{eq:cyclotomic-detkill}
\end{equation}
This argument works integrally at primes dividing \(d\).

Define the integers
\[
 D=Mk(k+1)!,\qquad S_0=2a^kM^k k!,\qquad B=D S_0,\qquad q=d/2 .
\]
The denominators of all \(B_i\) for \(0\le i\le k\) divide
\((k+1)!\) by von Staudt--Clausen, including \(B_1=-1/2\).
Expanding \(E_k\) thus shows \(D E_k(M,x)\in\mathbb Z\) for each
integer \(x\), since the smallest power of \(M\) in the expansion
is \(M^{-1}\). It follows that \(DY\in\mathbb Z[\Gamma]\), and
\(X=D^q\Delta\) is an integer.

We prove \(0<|X|\le B^q\).
Over \(\mathbb C\), the eigenvalues of multiplication by \(Y\) on
the quotient are its evaluations at the characters \(\lambda\) of
\(\Gamma\) with \(\lambda(c)=\epsilon\). This follows by the
finite Fourier transform on the abelian group \(\Gamma\).
Set \(\xi(b)=\lambda(\pi_M(b))\), extended by zero on nonunits.
Then
\[
 \lambda(Y)=(1-a^k\xi(a))\,L(1-k,\xi).
\]
For completeness, the identity for the second factor follows from
\(L(s,\xi)=M^{-s}\sum_{b\in G_M}\xi(b)\zeta(s,[b]_M/M)\) and
\(\zeta(1-k,x)=-B_k(x)/k\). Here \(\zeta(s,x)\) is the Hurwitz
zeta function. If \(\xi^*\) induces \(\xi\) and has conductor \(f\),
then
\[
 L(1-k,\xi)=L(1-k,\xi^*)
   \prod_{\substack{p\mid M\\p\nmid f}}(1-\xi^*(p)p^{k-1}).
\]
Each omitted Euler factor is nonzero as \(k\ge2\). The gamma factor
of the primitive completed \(L\)-function is finite and nonzero at
\(1-k\). Indeed, writing \(\xi^*(-1)=(-1)^\delta\), its argument is
\((1-k+\delta)/2\), an integer plus one half, since
\(\delta\equiv k\bmod 2\). The primitive functional equation
relates the completed value there to its value for
\(\overline{\xi^*}\) at \(k>1\). The latter is finite and nonzero
by the absolutely convergent Euler product. This includes the
primitive trivial character and the zeta functional equation.
Finally \(|a^k\xi(a)|=a^k>1\). Every eigenvalue is therefore nonzero.

The Fourier series of \(B_k\) gives, for \(0\le u\le1\) and \(k\ge2\),
\[
 |B_k(u)|\le \frac{2k!}{(2\pi)^k}\sum_{n\ge1}n^{-k}\le k!.
\]
The sum for \(\Theta(M)\) has at most \(M\) terms, while
\(|1-a^k\xi(a)|\le2a^k\). Hence \(|\lambda(Y)|\le S_0\).
The product of these \(q\) eigenvalues gives
\(0<|X|=D^q|\Delta|\le B^q\).

Now set \(r=v_\ell(|X|)+v_\ell(\varphi(M))\).
Since \(\Delta\in\mathbb Z_\ell\) and \(X=D^q\Delta\),
\(v_\ell(\Delta)\le v_\ell(|X|)\). By
\eqref{eq:cyclotomic-detkill}, \(\ell^r x=0\) for all
\(x\in P_\epsilon\). Furthermore,
\[
 r\log\ell \le q\log B+\log M
                  \le d\log(BM).
\]
For \(T=3e\ell N\), we have \(M,k,2\le T\), \(a\le T^3\), and
\(k+1\le T^2\). Bounding \(k!\) by \(T^k\) and \((k+1)!\) by
\(T^{2(k+1)}\) gives \(BM\le T^{7k+6}\le T^{10k}\) because
\(k\ge2\). Since \(e\ell N\ge2\), \(T\le(e\ell N)^3\).
This proves \eqref{eq:cyclotomic-kill}.

For the length assertion, multiplication by \(\ell^j\) embeds
\(H[\ell^{j+1}]/H[\ell^j]\) into \(H[\ell]\) for every \(j\ge0\).
Section~\ref{sec:selmer-generators} proves
\(\operatorname{length}_{\mathbb Z_\ell}H[\ell]\le
8\ell d\log(\ell N)\). Filtering \(P_\epsilon\subset H[\ell^r]\)
by the successive kernels and applying \eqref{eq:cyclotomic-kill}
gives \eqref{eq:cyclotomic-length}.

\paragraph{Correspondence with the release.}
In the \texttt{OddResidueBounds} development, the class relation is
\texttt{clKill} and its reciprocity consequence is \texttt{homKill}.
The Bernoulli approximation is \texttt{congr} in \texttt{NearLevel}.
The restriction, power and integral annihilators are
\texttt{resKill}, \texttt{totKill} and \texttt{zKill}.
The finite-layer argument above uses the original Selmer cocycle
and its finitely many local trivializations to construct the
unramified character directly. The release obtains the same character
on a sufficiently deep layer by descent of its restriction.
The determinant, logarithmic estimate, exponent and filtration
correspond to \texttt{detVal}, \texttt{bdxLog}, \texttt{killPow}
and \texttt{lenPow}, respectively. The conclusion is the source's
\texttt{StickBound} with \(C=240\).
